\textbf{COLUNI 2019 -- Questão 10}\par\medskip No filme \textit{O carteiro e o poeta}, o personagem Mário e o poeta Pablo Neruda compreendem que a metáfora é o elemento constituinte da poesia. A poesia de Manoel de Barros também apresenta a metáfora como elemento constituinte.

\par

Assinale a alternativa em que o excerto de Manoel de Barros explica a condição para o uso da metáfora em sua poesia:\begin{enumerate}
\item “Eu queria então escovar as palavras para escutar o primeiro esgar de cada uma.” (BARROS, M. \textbf{Memórias Inventadas}. Rio
      de Janeiro: Alfaguara, 2018, p.17)
\item “[...] quando um sujeito sabe quase tudo, ele não faz a imagem, ele solta a palavra [...]”. (BARROS, M. \textbf{Memórias
      Inventadas}. Rio de Janeiro: Alfaguara, 2018, p. 9)
\item “É uma indigência vocabular que provoca a imagem.” (BARROS, M. \textbf{Memórias Inventadas}. Rio de Janeiro: Alfaguara, 2018, p. 9)
\item “Descobri nesse tempo que os apelidos pregam mais quando trovam [...]”. (BARROS, M. \textbf{Memórias Inventadas}. Rio de
      Janeiro: Alfaguara, 2018, p. 45)

\end{enumerate}
